% Ilustrasi rentang nilai entropi Shannon: 0 <= H(X) <= log2(n).
% Dirujuk dari section-02.tex dengan % Ilustrasi rentang nilai entropi Shannon: 0 <= H(X) <= log2(n).
% Dirujuk dari section-02.tex dengan % Ilustrasi rentang nilai entropi Shannon: 0 <= H(X) <= log2(n).
% Dirujuk dari section-02.tex dengan % Ilustrasi rentang nilai entropi Shannon: 0 <= H(X) <= log2(n).
% Dirujuk dari section-02.tex dengan \input{figures/rentang-entropi}.
\begin{figure}[htbp]
    \centering
    \begin{tikzpicture}[font=\small]
        % Batang rentang entropi, diwarnai bertingkat dari terang ke gelap
        \foreach \i in {0,...,5} {
            \pgfmathsetmacro{\xa}{\i*10/6}
            \pgfmathsetmacro{\xb}{(\i+1)*10/6}
            \fill[blue!\the\numexpr 8+11*\i\relax]
                (\xa,0.35) rectangle (\xb,1.00);
        }
        \draw[blue!70!black, thick] (0,0.35) rectangle (10,1.00);

        % Sumbu nilai entropi
        \draw[-{Latex[length=2.4mm]}, semithick] (0,0.05) -- (10.9,0.05);
        \node[below, font=\scriptsize, inner sep=1.5pt] at (0.25,-0.02) {$0$};
        \node[below, font=\scriptsize, inner sep=1.5pt] at (10.4,-0.02)
            {$\log_2 n$};

        % Anotasi kedua ujung rentang
        \node[above, align=center, font=\scriptsize] at (0.9,1.10)
            {$H = 0$\\ satu simbol berpeluang $1$};
        \node[above, align=center, font=\scriptsize] at (9.2,1.10)
            {$H = \log_2 n$\\ semua simbol berpeluang $1/n$};

        % Garis penunjuk ke keterangan di bawah
        \draw[gray!70, densely dotted] (0,0.35) -- (0,-0.32);
        \draw[gray!70, densely dotted] (10,0.35) -- (10,-0.32);

        \node[below, align=center, font=\scriptsize, text width=4.0cm]
            at (1.9,-0.32)
            {Pesan sudah pasti: tidak ada ketidakpastian yang tersisa};
        \node[below, align=center, font=\scriptsize, text width=4.0cm]
            at (8.3,-0.32)
            {Ketidakpastian paling tinggi: tiap simbol membawa informasi maksimal};
    \end{tikzpicture}
    \caption{Rentang nilai entropi Shannon $0 \le H(X) \le \log_2 n$. Nilai minimum tercapai bila pesan sudah pasti, nilai maksimum bila seluruh $n$ simbol muncul dengan peluang sama}
    \label{fig:rentang-entropi}
\end{figure}
.
\begin{figure}[htbp]
    \centering
    \begin{tikzpicture}[font=\small]
        % Batang rentang entropi, diwarnai bertingkat dari terang ke gelap
        \foreach \i in {0,...,5} {
            \pgfmathsetmacro{\xa}{\i*10/6}
            \pgfmathsetmacro{\xb}{(\i+1)*10/6}
            \fill[blue!\the\numexpr 8+11*\i\relax]
                (\xa,0.35) rectangle (\xb,1.00);
        }
        \draw[blue!70!black, thick] (0,0.35) rectangle (10,1.00);

        % Sumbu nilai entropi
        \draw[-{Latex[length=2.4mm]}, semithick] (0,0.05) -- (10.9,0.05);
        \node[below, font=\scriptsize, inner sep=1.5pt] at (0.25,-0.02) {$0$};
        \node[below, font=\scriptsize, inner sep=1.5pt] at (10.4,-0.02)
            {$\log_2 n$};

        % Anotasi kedua ujung rentang
        \node[above, align=center, font=\scriptsize] at (0.9,1.10)
            {$H = 0$\\ satu simbol berpeluang $1$};
        \node[above, align=center, font=\scriptsize] at (9.2,1.10)
            {$H = \log_2 n$\\ semua simbol berpeluang $1/n$};

        % Garis penunjuk ke keterangan di bawah
        \draw[gray!70, densely dotted] (0,0.35) -- (0,-0.32);
        \draw[gray!70, densely dotted] (10,0.35) -- (10,-0.32);

        \node[below, align=center, font=\scriptsize, text width=4.0cm]
            at (1.9,-0.32)
            {Pesan sudah pasti: tidak ada ketidakpastian yang tersisa};
        \node[below, align=center, font=\scriptsize, text width=4.0cm]
            at (8.3,-0.32)
            {Ketidakpastian paling tinggi: tiap simbol membawa informasi maksimal};
    \end{tikzpicture}
    \caption{Rentang nilai entropi Shannon $0 \le H(X) \le \log_2 n$. Nilai minimum tercapai bila pesan sudah pasti, nilai maksimum bila seluruh $n$ simbol muncul dengan peluang sama}
    \label{fig:rentang-entropi}
\end{figure}
.
\begin{figure}[htbp]
    \centering
    \begin{tikzpicture}[font=\small]
        % Batang rentang entropi, diwarnai bertingkat dari terang ke gelap
        \foreach \i in {0,...,5} {
            \pgfmathsetmacro{\xa}{\i*10/6}
            \pgfmathsetmacro{\xb}{(\i+1)*10/6}
            \fill[blue!\the\numexpr 8+11*\i\relax]
                (\xa,0.35) rectangle (\xb,1.00);
        }
        \draw[blue!70!black, thick] (0,0.35) rectangle (10,1.00);

        % Sumbu nilai entropi
        \draw[-{Latex[length=2.4mm]}, semithick] (0,0.05) -- (10.9,0.05);
        \node[below, font=\scriptsize, inner sep=1.5pt] at (0.25,-0.02) {$0$};
        \node[below, font=\scriptsize, inner sep=1.5pt] at (10.4,-0.02)
            {$\log_2 n$};

        % Anotasi kedua ujung rentang
        \node[above, align=center, font=\scriptsize] at (0.9,1.10)
            {$H = 0$\\ satu simbol berpeluang $1$};
        \node[above, align=center, font=\scriptsize] at (9.2,1.10)
            {$H = \log_2 n$\\ semua simbol berpeluang $1/n$};

        % Garis penunjuk ke keterangan di bawah
        \draw[gray!70, densely dotted] (0,0.35) -- (0,-0.32);
        \draw[gray!70, densely dotted] (10,0.35) -- (10,-0.32);

        \node[below, align=center, font=\scriptsize, text width=4.0cm]
            at (1.9,-0.32)
            {Pesan sudah pasti: tidak ada ketidakpastian yang tersisa};
        \node[below, align=center, font=\scriptsize, text width=4.0cm]
            at (8.3,-0.32)
            {Ketidakpastian paling tinggi: tiap simbol membawa informasi maksimal};
    \end{tikzpicture}
    \caption{Rentang nilai entropi Shannon $0 \le H(X) \le \log_2 n$. Nilai minimum tercapai bila pesan sudah pasti, nilai maksimum bila seluruh $n$ simbol muncul dengan peluang sama}
    \label{fig:rentang-entropi}
\end{figure}
.
\begin{figure}[htbp]
    \centering
    \begin{tikzpicture}[font=\small]
        % Batang rentang entropi, diwarnai bertingkat dari terang ke gelap
        \foreach \i in {0,...,5} {
            \pgfmathsetmacro{\xa}{\i*10/6}
            \pgfmathsetmacro{\xb}{(\i+1)*10/6}
            \fill[blue!\the\numexpr 8+11*\i\relax]
                (\xa,0.35) rectangle (\xb,1.00);
        }
        \draw[blue!70!black, thick] (0,0.35) rectangle (10,1.00);

        % Sumbu nilai entropi
        \draw[-{Latex[length=2.4mm]}, semithick] (0,0.05) -- (10.9,0.05);
        \node[below, font=\scriptsize, inner sep=1.5pt] at (0.25,-0.02) {$0$};
        \node[below, font=\scriptsize, inner sep=1.5pt] at (10.4,-0.02)
            {$\log_2 n$};

        % Anotasi kedua ujung rentang
        \node[above, align=center, font=\scriptsize] at (0.9,1.10)
            {$H = 0$\\ satu simbol berpeluang $1$};
        \node[above, align=center, font=\scriptsize] at (9.2,1.10)
            {$H = \log_2 n$\\ semua simbol berpeluang $1/n$};

        % Garis penunjuk ke keterangan di bawah
        \draw[gray!70, densely dotted] (0,0.35) -- (0,-0.32);
        \draw[gray!70, densely dotted] (10,0.35) -- (10,-0.32);

        \node[below, align=center, font=\scriptsize, text width=4.0cm]
            at (1.9,-0.32)
            {Pesan sudah pasti: tidak ada ketidakpastian yang tersisa};
        \node[below, align=center, font=\scriptsize, text width=4.0cm]
            at (8.3,-0.32)
            {Ketidakpastian paling tinggi: tiap simbol membawa informasi maksimal};
    \end{tikzpicture}
    \caption{Rentang nilai entropi Shannon $0 \le H(X) \le \log_2 n$. Nilai minimum tercapai bila pesan sudah pasti, nilai maksimum bila seluruh $n$ simbol muncul dengan peluang sama}
    \label{fig:rentang-entropi}
\end{figure}
